\documentclass{article} % For LaTeX2e
\usepackage{iclr_templates/iclr2025_conference,times}
\iclrfinalcopy
% Optional math commands from https://github.com/goodfeli/dlbook_notation.
\input{iclr_templates/math_commands}

\usepackage{hyperref}
\usepackage{url}
\usepackage{microtype}
\usepackage{graphicx}
\usepackage{enumitem}
\usepackage{subfigure}
\usepackage{booktabs} % for professional tables
\usepackage{wrapfig}
\usepackage{booktabs}
\usepackage{multirow}
\usepackage{adjustbox} % optional, for resizing large tables
\usepackage{xcolor,colortbl}
\usepackage{pgfmath}
\usepackage{rotating}
\usepackage{siunitx}
\usepackage{makecell}
\usepackage{algorithm}
\usepackage{algpseudocode}
% \usepackage[table]{xcolor}
\input{color_commands}
\sisetup{
  detect-weight=true,
  detect-inline-weight=math,
  table-format=1.3,
  table-number-alignment=center,
}



\newcommand{\thewell}{The Well}
\newcommand{\acoustic}{\texttt{acoustic\_scattering}}
\newcommand{\activematter}{\texttt{active\_matter}}
\newcommand{\rsg}{\texttt{convective\_envelope\_rsg}}
\newcommand{\euler}{\texttt{euler\_multi\_quadrants}}
\newcommand{\staircase}{\texttt{helmholtz\_staircase}}
\newcommand{\MHD}{\texttt{MHD}}
\newcommand{\MHDsmall}{\texttt{MHD\_64}}
\newcommand{\MHDbig}{\texttt{MHD\_256}}
\newcommand{\pattern}{\texttt{gray\_scott\_reaction\_diffusion}}
\newcommand{\planetswe}{\texttt{planetswe}}
\newcommand{\neutronstar}{\texttt{post\_neutron\_star\_merger}}
\newcommand{\RB}{\texttt{rayleigh\_benard}}
\newcommand{\RT}{\texttt{rayleigh\_taylor\_instability}}
\newcommand{\shearflow}{\texttt{shear\_flow}}
\newcommand{\supernova}{\texttt{supernova\_explosion}}
\newcommand{\TGC}{\texttt{turbulence\_gravity\_cooling}}
\newcommand{\TRLtwo}{\texttt{turbulent\_radiative\_layer\_2D}}
\newcommand{\TRLthree}{\texttt{turbulent\_radiative\_layer\_3D}}
\newcommand{\viscoelastic}{\texttt{viscoelastic\_instability}}
\DeclareSymbolFont{matha}{OML}{txmi}{m}{it} % txfonts
\DeclareMathSymbol{\varv}{\mathord}{matha}{118}
\newcommand{\cartesiantwoD}{($x,y$)}
\newcommand{\cartesianthreeD}{($x,y,z$)}
\newcommand{\spherical}{($r,\theta,\phi$)}
\newcommand{\logspherical}{($\log r,\theta,\phi$)}
\newcommand{\angular}{($\theta,\phi$)}
\newcommand{\VPP}{\textbf{VPP: }}
\newcommand{\Fields}{\textbf{Fields: }}
\newcommand{\Chall}{\textbf{Challenges: }}
\renewcommand{\vec}[1]{\boldsymbol{#1}}
\newcommand\grad{\boldsymbol{\nabla}}    
\newcommand{\best}[1]{\cellcolor{yellow!30}#1}
\title{Making LodeRunner into a diffusion model.}


% Authors must not appear in the submitted version. They should be hidden
% as long as the \iclrfinalcopy macro remains commented out below.
% Non-anonymous submissions will be rejected without review.
% \newcommand{\mysep}{\hspace{12pt}}

 
% The \author macro works with any number of authors. There are two commands
% used to separate the names and addresses of multiple authors: \And and \AND.
%
% Using \And between authors leaves it to \LaTeX{} to determine where to break
% the lines. Using \AND forces a linebreak at that point. So, if \LaTeX{}
% puts 3 of 4 authors names on the first line, and the last on the second
% line, try using \AND instead of \And before the third author name.
\author{Kyle Hickmann\thanks{Contact: \texttt{hickmank@lanl.gov}}$\;^{1}$,
   \vspace{8pt}
   \\
   \textnormal{$^{1\,}$Los Alamos National Laboratory},   
}

\newcommand{\fix}{\marginpar{FIX}}
\newcommand{\new}{\marginpar{NEW}}
\newcommand{\Walrus}{\texttt{Walrus}}

\newcommand{\shadecell}[2]{%
  % #1 = value (0–1), #2 = text to show
  \pgfmathsetmacro{\intensity}{100*(1-#1)} % color intensity
  \ifdim #1pt>0.6pt
    \cellcolor{orange!\intensity!white}{\textbf{\textcolor{orange!80!black}{#2}}}%
  \else
    \cellcolor{orange!\intensity!white}{\textcolor{orange!80!black}{#2}}%
  \fi
}
%\iclrfinalcopy % Uncomment for camera-ready version, but NOT for submission.
\begin{document}


\maketitle



\begin{abstract}
Plan to create a diffusion model from LodeRuner architecture.
\end{abstract}

\section{LodeRunner: mathematical description}

% bomberman.py — Mathematical description of the LodeRunner architecture (journal style)

\subsection{LodeRunner: Variable-Conditioned Patch-Token Surrogate with SWIN U-Net Core}
\label{subsec:loderunner_math}

We consider a variable-conditioned surrogate model that maps a collection of spatial fields to another collection of spatial fields at a specified lead time. The \texttt{LodeRunner} architecture (in \texttt{src/yoke/models/vit/swin/bomberman.py}) implements this mapping by (i) representing the input as a set of patch tokens, (ii) fusing information across physical variables through attention, (iii) applying a hierarchical windowed self-attention U-Net (Swin U-Net) on the token grid, and (iv) decoding back to full-resolution fields.

\paragraph{Problem setting.}
Let $\Omega = \{1,\dots,H\}\times\{1,\dots,W\}$ be a discrete spatial grid. Let $\mathcal{V}$ be a finite vocabulary of variables (e.g., pressure, temperature), with $|\mathcal{V}| = V$. For each sample, we observe a subset of input variables $S_{\mathrm{in}}\subseteq \mathcal{V}$ and wish to predict a subset of output variables $S_{\mathrm{out}}\subseteq \mathcal{V}$ at lead time $\tau$.
The input is represented as a multi-channel image
\[
x \in \mathbb{R}^{|S_{\mathrm{in}}|\times H\times W},
\]
together with index sets $S_{\mathrm{in}}$ and $S_{\mathrm{out}}$, and a lead-time scalar (or vector) $\tau$. The model outputs
\[
\hat{y} \in \mathbb{R}^{|S_{\mathrm{out}}|\times H\times W}.
\]
In \texttt{bomberman.py}, the variables are indexed relative to a fixed list \texttt{default\_vars} of size $V$; the network forms predictions for all $V$ variables and then selects the requested subset $S_{\mathrm{out}}$.

\subsubsection{Patch-token representation}
\label{subsec:patch_token_math}
Fix a patch size $(P_h,P_w)$ and partition $\Omega$ into a regular grid of $N$ non-overlapping patches $\{ \Omega_n \}_{n=1}^N$ (so $N=(H/P_h)(W/P_w)$ under exact divisibility). For each variable $v\in \mathcal{V}$, define a patchification operator
\[
\Pi(x_v) \in \mathbb{R}^{N \times (P_hP_w)},
\]
which collects the pixels of $x_v$ inside each patch into a length-$P_hP_w$ vector. A learned embedding map then lifts each patch vector to a $D$-dimensional token:
\[
E:\mathbb{R}^{P_hP_w}\to \mathbb{R}^D.
\]
Applied in parallel across variables present in $S_{\mathrm{in}}$, this yields an initial set of tokens
\[
z^{(0)}_{v,n} = E(\Pi(x_v)_n),\qquad v\in S_{\mathrm{in}},\; n=1,\dots,N.
\]
This step corresponds to the ``parallel'' patch embedding in \texttt{bomberman.py}.

\subsubsection{Variable identity encoding and cross-variable fusion}
\label{subsec:var_fusion_math}
Because each token arises from a specific physical variable, the architecture injects a learnable variable identifier. Let $e_v\in\mathbb{R}^D$ denote a learned embedding for variable $v\in\mathcal{V}$. The variable-tagged tokens are
\[
z^{(1)}_{v,n} = z^{(0)}_{v,n} + e_v.
\]
To combine information across variables at the same spatial patch index, LodeRunner applies an attention-based aggregation over the variable axis. For each patch index $n$, define
\[
z^{(2)}_n = \mathrm{Agg}\big(\{z^{(1)}_{v,n}\}_{v\in S_{\mathrm{in}}}\big)\in\mathbb{R}^D,
\]
where $\mathrm{Agg}$ is a multi-head attention mechanism (parameterized by \texttt{num\_heads} in \texttt{bomberman.py}). Intuitively, $\mathrm{Agg}$ produces a single fused token per patch that summarizes the multi-variable state while allowing the model to weight variables adaptively (context-dependent coupling).

\subsubsection{Spatial and temporal conditioning}
\label{subsec:pos_time_math}
The fused patch tokens are indexed by the patch grid location. LodeRunner adds a learned positional encoding $p_n\in\mathbb{R}^D$:
\[
z^{(3)}_n = z^{(2)}_n + p_n,
\]
which provides discrete spatial reference at the patch level (``index-aware'' positional information, per \texttt{bomberman.py}).

To condition on lead time, a learned time embedding $t(\tau)\in\mathbb{R}^D$ is injected:
\[
z^{(4)}_n = z^{(3)}_n + t(\tau).
\]
This allows the network to represent a family of time-dependent mappings using shared parameters, with $\tau$ entering as a global conditioning signal.

\subsubsection{Hierarchical windowed-attention U-Net on the token grid}
\label{subsec:swin_unet_math}
Let $Z^{(4)}\in \mathbb{R}^{N\times D}$ collect the tokens $\{z^{(4)}_n\}_{n=1}^N$ arranged on the 2D patch grid. The core operator is a hierarchical encoder--decoder
\[
Z^{(5)} = \mathcal{U}(Z^{(4)}),
\]
where $\mathcal{U}$ is a U-Net-like composition of (i) local windowed self-attention blocks and (ii) multi-scale patch merging/expansion. The locality is governed by window sizes (given by \texttt{window\_sizes} in \texttt{bomberman.py}), and the change of resolution is controlled by fixed merge scales (given by \texttt{patch\_merge\_scales}). The channel dimension varies across levels by a multiplicative factor (\texttt{emb\_factor}). The number of attention blocks per level is set by \texttt{block\_structure}. This design yields (a) tractable attention complexity at high resolution via windowing, and (b) long-range effective receptive fields via the hierarchical multi-scale pathway, analogous to classical U-Nets.

\subsubsection{Token-to-field decoding via linear patch reconstruction}
\label{subsec:decode_math}
From the final token representation, the model reconstructs per-variable patch contents via a linear map
\[
\phi:\mathbb{R}^D \to \mathbb{R}^{V\cdot P_hP_w}.
\]
For each patch index $n$,
\[
r_n = \phi(z^{(5)}_n) \in \mathbb{R}^{V\cdot P_hP_w},
\]
which can be reshaped into $\{r_{v,n}\in\mathbb{R}^{P_hP_w}\}_{v\in\mathcal{V}}$. An ``unpatchify'' operator $\Pi^{-1}$ then places each patch vector back onto the spatial grid to obtain full-resolution predictions for all variables:
\[
\hat{Y} = \Pi^{-1}(\{r_{v,n}\}_{v\in\mathcal{V},\,n=1}^N) \in \mathbb{R}^{V\times H\times W}.
\]
Finally, the requested outputs are selected:
\[
\hat{y} = \hat{Y}\big|_{S_{\mathrm{out}}} \in \mathbb{R}^{|S_{\mathrm{out}}|\times H\times W}.
\]

\subsubsection{Noise-regularized input (optional)}
\label{subsec:noise_math}
An optional noise injection is applied to the input fields prior to embedding:
\[
\tilde{x} = x + \sigma \,\|x\|_2\,\varepsilon,\qquad \varepsilon\sim\mathcal{N}(0,I),
\]
with scalar $\sigma\ge 0$ (the \texttt{noise\_scale} parameter in \texttt{bomberman.py}), computed per sample. This acts as a scale-aware perturbation that may improve robustness and generalization.

\paragraph{Summary.}
In functional form, the architecture defines a mapping
\[
F_\theta:\ (x, S_{\mathrm{in}}, S_{\mathrm{out}}, \tau)\ \mapsto\ \hat{y},
\]
implemented by patch-tokenization and variable tagging, attention-based fusion across variables, addition of spatial and temporal embeddings, a hierarchical windowed-attention U-Net $\mathcal{U}$ operating on the patch grid, and a linear patch decoder followed by unpatchification and output selection. This yields a flexible multi-variable, lead-time-conditioned surrogate suitable for multi-physics settings while preserving an image-to-image inductive bias through the patch grid structure.


\section{Extension of LodeRunner to Score-based Diffusion}

% loderunner_diffusion.tex — Mathematical description of a score-based diffusion extension of LodeRunner (journal style)

\subsection{Diffusion LodeRunner: Conditional Score Model for Probabilistic Forecasting}
\label{subsec:loderunner_diffusion_math}

We extend the deterministic \texttt{LodeRunner} mapping in \texttt{src/yoke/models/vit/swin/bomberman.py}
to a \emph{conditional score-based diffusion model} that represents a full conditional distribution
over future fields rather than a single point estimate. The resulting model, which we refer to as
\texttt{DiffusionLodeRunner}, learns to denoise a stochastically perturbed version of the target fields
while conditioning on the observed input fields and lead time.

\paragraph{Problem setting.}
Let $x\in\mathbb{R}^{|S_{\mathrm{in}}|\times H\times W}$ denote the input multi-field state at an initial time,
with variable index set $S_{\mathrm{in}}\subseteq\mathcal{V}$, and let $y\in\mathbb{R}^{|S_{\mathrm{out}}|\times H\times W}$
denote the desired output state at lead time $\tau$ (time offset), with $S_{\mathrm{out}}\subseteq\mathcal{V}$.
In the temporal LSC dataset (\texttt{src/yoke/datasets/lsc\_dataset.py}), $(x,y,\tau)$ corresponds to
$(\texttt{start\_img},\texttt{end\_img},\texttt{Dt})$.

The deterministic surrogate previously implemented by \texttt{LodeRunner} is a conditional map
\[
F_\theta:\ (x,S_{\mathrm{in}},S_{\mathrm{out}},\tau)\ \mapsto\ \hat{y}.
\]
We replace this with a conditional diffusion model that enables sampling:
\[
p_\theta(y\mid x,S_{\mathrm{in}},S_{\mathrm{out}},\tau).
\]

\subsubsection{Variance-preserving (VP) forward diffusion on the target}
\label{subsec:vp_forward_diffusion}

We introduce a diffusion time parameter $t\in[0,1]$ and define a \emph{variance-preserving} (VP) forward
noising process applied to the target fields $y$. Let $\epsilon\sim\mathcal{N}(0,I)$ be i.i.d. Gaussian noise
with the same shape as $y$. We define the perturbed target at diffusion time $t$ as
\begin{equation}
y_t \;=\; \alpha(t)\,y \;+\; \sigma(t)\,\epsilon,
\label{eq:vp_forward}
\end{equation}
where $\alpha(t),\sigma(t)\in\mathbb{R}$ are scalar noise schedule functions satisfying
\begin{equation}
\alpha(t)^2 + \sigma(t)^2 = 1,\qquad \alpha(0)=1,\ \sigma(0)=0,\qquad \alpha(1)\approx 0,\ \sigma(1)\approx 1.
\label{eq:vp_constraint}
\end{equation}
This construction preserves the marginal variance of $y_t$ across diffusion time and provides a continuous
family of corruption levels.

\subsubsection{Score parameterization via noise prediction}
\label{subsec:score_param}

Score-based diffusion modeling seeks the score
\[
s(y_t,t\mid x,S_{\mathrm{in}},S_{\mathrm{out}},\tau) \;:=\; \nabla_{y_t}\log p(y_t\mid x,S_{\mathrm{in}},S_{\mathrm{out}},\tau,t).
\]
Rather than predicting the score directly, we parameterize it through noise prediction, following standard
denoising diffusion practice. We define a neural network
\begin{equation}
\epsilon_\theta:\ (x,y_t,S_{\mathrm{in}},S_{\mathrm{out}},\tau,t)\ \mapsto\ \widehat{\epsilon},
\label{eq:eps_model}
\end{equation}
and recover an implied score estimate via the VP identity
\begin{equation}
\widehat{s}_\theta(y_t,t\mid x,\tau) \;\approx\; -\frac{1}{\sigma(t)}\,\widehat{\epsilon}.
\label{eq:score_from_eps}
\end{equation}
Equivalently, an estimate of the clean target is obtained by
\begin{equation}
\widehat{y}_0 \;=\; \frac{y_t - \sigma(t)\widehat{\epsilon}}{\alpha(t)}.
\label{eq:y0_hat}
\end{equation}

\subsubsection{Architecture modifications: dual-stream tokenization and diffusion-time conditioning}
\label{subsec:arch_mods}

The original \texttt{LodeRunner} architecture first tokenizes the input fields $x$ via parallel patch embedding,
injects variable tags, aggregates across variables, adds positional and lead-time conditioning, and then
applies a SWIN U-Net operator on the token grid (cf.\ \texttt{src/yoke/loderunner.tex}).

To condition denoising on both the observed state $x$ and the corrupted target $y_t$ while \emph{preserving the
variable-subset flexibility}, we introduce a \emph{dual-stream} tokenization and fusion mechanism:

\paragraph{Stream A (conditioning stream).}
Using the existing parallel patch embedding and variable fusion pipeline, we produce fused patch tokens
from the conditioning variables $S_{\mathrm{in}}$:
\[
Z^{x} \;=\; \mathrm{FuseTokens}\big(x,S_{\mathrm{in}}\big)\ \in\ \mathbb{R}^{N\times D}.
\]
Here $\mathrm{FuseTokens}(\cdot)$ denotes the composition of parallel patch embedding, variable identity
encoding, attention-based aggregation over variables, and positional encoding, as defined in
Sections~\ref{subsec:patch_token_math}--\ref{subsec:pos_time_math} of \texttt{src/yoke/loderunner.tex}.

\paragraph{Stream B (noised-target stream).}
Analogously, we tokenize the noised target variables in $S_{\mathrm{out}}$:
\[
Z^{y_t} \;=\; \mathrm{FuseTokens}\big(y_t,S_{\mathrm{out}}\big)\ \in\ \mathbb{R}^{N\times D}.
\]

\paragraph{Token fusion.}
Because both streams are aggregated to one token per patch index, they share the same token grid resolution.
We fuse them by addition:
\begin{equation}
Z^{(0)} \;=\; Z^{x} + Z^{y_t}\ \in\ \mathbb{R}^{N\times D}.
\label{eq:token_fusion}
\end{equation}
(Other fusion operators are possible, e.g.\ concatenation with a learned projection, but the additive form
maintains dimensional compatibility with the existing SWIN U-Net backbone.)

\paragraph{Lead-time and diffusion-time conditioning.}
The original lead-time embedding $t(\tau)\in\mathbb{R}^D$ is retained. We additionally introduce a diffusion-time
embedding $g(t)\in\mathbb{R}^D$ of the same functional form (a learned linear map of the scalar time):
\begin{equation}
Z^{(1)}_n \;=\; Z^{(0)}_n \;+\; t(\tau) \;+\; g(t),\qquad n=1,\dots,N.
\label{eq:time_conditioning}
\end{equation}
This mirrors the implementation pattern of \texttt{TimeEmbed} in \texttt{src/yoke/models/vit/embedding\_encoders.py},
applied once for $\tau$ and once for $t$.

\paragraph{SWIN U-Net core and decoding.}
We reuse the hierarchical windowed-attention U-Net operator $\mathcal{U}$ and the token-to-field decoding
pipeline (linear patch reconstruction + unpatchification) unchanged:
\[
Z^{(2)} = \mathcal{U}(Z^{(1)}),\qquad
\widehat{E} = \Pi^{-1}\circ \phi(Z^{(2)}).
\]
Finally, we select the output variables and interpret them as a noise prediction:
\begin{equation}
\widehat{\epsilon} \;=\; \widehat{E}\big|_{S_{\mathrm{out}}}\ \in\ \mathbb{R}^{|S_{\mathrm{out}}|\times H\times W}.
\label{eq:eps_select}
\end{equation}
In contrast to the deterministic model, the decoded fields are not interpreted as $\hat{y}$ directly, but as
$\widehat{\epsilon}$ (or equivalently as a score via \eqref{eq:score_from_eps}).

\subsubsection{Training objective (denoising score matching)}
\label{subsec:training_objective}

Training proceeds by sampling $(x,y,\tau)$ from the dataset and sampling diffusion times $t\sim\mathrm{Uniform}(0,1)$
independently per batch element. We generate $y_t$ by \eqref{eq:vp_forward} and minimize the mean-squared error
between predicted and true noise:
\begin{equation}
\mathcal{L}(\theta)
\;=\;
\mathbb{E}_{(x,y,\tau)}\,
\mathbb{E}_{t\sim \mathcal{U}(0,1)}\,
\mathbb{E}_{\epsilon\sim\mathcal{N}(0,I)}
\left[
\left\|
\epsilon_\theta(x,y_t,S_{\mathrm{in}},S_{\mathrm{out}},\tau,t)
-
\epsilon
\right\|_2^2
\right].
\label{eq:diffusion_loss}
\end{equation}
This objective is equivalent to denoising score matching under the VP perturbation kernel and induces the
conditional score model \eqref{eq:score_from_eps}.

\subsubsection{Sampling (DDIM-style probability flow approximation)}
\label{subsec:sampling_ddim}

Given a conditioning input $(x,S_{\mathrm{in}},S_{\mathrm{out}},\tau)$, we sample from the learned conditional
distribution by initializing $y_{t_N}\sim\mathcal{N}(0,I)$ on the output variables and iterating a decreasing
sequence of diffusion times $1=t_N > t_{N-1} > \dots > t_0\approx 0$. At each step, we compute
\[
\widehat{\epsilon}_i = \epsilon_\theta(x,y_{t_i},S_{\mathrm{in}},S_{\mathrm{out}},\tau,t_i),
\qquad
\widehat{y}_0 = \frac{y_{t_i}-\sigma(t_i)\widehat{\epsilon}_i}{\alpha(t_i)},
\]
and update deterministically (DDIM with $\eta=0$):
\begin{equation}
y_{t_{i-1}}
\;=\;
\alpha(t_{i-1})\,\widehat{y}_0 \;+\; \sigma(t_{i-1})\,\widehat{\epsilon}_i.
\label{eq:ddim_update}
\end{equation}
The final sample $y_{t_0}$ (or $\widehat{y}_0$ at the final step) constitutes a probabilistic forecast of the
requested output fields.

\paragraph{Remark (noise injection in the deterministic model).}
The optional input noise injection described in Section~\ref{subsec:noise_math} of \texttt{src/yoke/loderunner.tex}
is not used during diffusion training, as the forward corruption \eqref{eq:vp_forward} defines the intended
noise distribution and scale for the denoising objective.

\bibliography{main}
\bibliographystyle{iclr_templates/iclr2025_conference}

% \newpage
\appendix
% \section{Appendix}

%\include{supplements}

\section{Supplemental: LodeRunner PyTorch implementation}

% bomberman.py — LaTeX description of the LodeRunner architecture (under ~2 pages)

\subsection{LodeRunner Architecture (Parallel Variable-Patch SWIN U-Net)}
\label{subsec:loderunner}

The \texttt{LodeRunner} module (defined in \texttt{src/yoke/models/vit/swin/bomberman.py}) is a variable-channel image-to-image network that maps an input tensor of spatial fields to an output tensor of spatial fields, while explicitly conditioning on (i) which variables are present in the input, (ii) which variables are requested as outputs, and (iii) a lead-time (temporal offset). Internally, it follows a \emph{patch tokenization} $\rightarrow$ \emph{variable-aware encoding/aggregation} $\rightarrow$ \emph{SWIN U-Net backbone} $\rightarrow$ \emph{unpatchification} pipeline.

\paragraph{Inputs and outputs.}
The forward signature is
\[
\texttt{forward}(x,\; \texttt{in\_vars},\; \texttt{out\_vars},\; \texttt{lead\_times}) \rightarrow \texttt{preds},
\]
where $x \in \mathbb{R}^{B \times C \times H \times W}$ is a batch of images/fields with $C$ variable channels, \texttt{in\_vars} specifies which variables are present in $x$, \texttt{out\_vars} specifies which variables to return, and \texttt{lead\_times} provides a temporal conditioning signal. The network returns
\[
\texttt{preds} = \texttt{preds}[:, \texttt{out\_vars}],
\]
i.e., it predicts all modeled variables and then selects the requested outputs (\texttt{bomberman.py}).

\paragraph{High-level computation graph.}
Let $P_h \times P_w$ be the patch size (\texttt{patch\_size}), and let $V$ be the maximum number of variables (\texttt{max\_vars = len(default\_vars)}). The main stages are:
\begin{enumerate}
    \item \textbf{Noise injection (optional regularization).}
    \item \textbf{Parallel patch embedding per variable channel.}
    \item \textbf{Variable-ID embedding (learnable variable tag).}
    \item \textbf{Variable aggregation via attention.}
    \item \textbf{Patch position encoding.}
    \item \textbf{Lead-time (temporal) encoding.}
    \item \textbf{SWIN U-Net backbone over patch tokens.}
    \item \textbf{Linear projection to per-variable patch pixels.}
    \item \textbf{Unpatchify to reconstruct full-resolution fields; select outputs.}
\end{enumerate}

\subsubsection{(1) Noise Injection}
\label{subsec:noise}
Before any embedding, \texttt{LodeRunner} applies scale-controlled Gaussian noise:
\[
\ell_2(x) = \sqrt{\sum_{c,h,w} x^2}\quad (\text{per-sample, keepdim}), \qquad
\tilde{x} = x + \sigma \,\ell_2(x)\,\epsilon,
\]
where $\epsilon \sim \mathcal{N}(0, I)$ and $\sigma=$ \texttt{noise\_scale} (\texttt{bomberman.py}). This implements magnitude-proportional perturbations that can improve robustness.

\subsubsection{(2) Parallel Variable Patch Embedding}
\label{subsec:parallel_patch_embed}
The first learnable stage is \texttt{ParallelVarPatchEmbed}, configured with:
\[
(H, W)=\texttt{image\_size},\quad (P_h,P_w)=\texttt{patch\_size},\quad D=\texttt{embed\_dim},\quad V=\texttt{max\_vars}.
\]
It converts the (possibly variable-channel) image into a sequence of patch tokens. The code comment states: ``Each channel is embedded independently'' (\texttt{bomberman.py}), implying a per-variable patchification/embedding into a common token dimension $D$. Denote the resulting patch-token representation by
\[
z_0 = \texttt{ParallelVarPatchEmbed}(\tilde{x}, \texttt{in\_vars}).
\]

\subsubsection{(3) Variable Embedding (Learnable Variable Tag)}
\label{subsec:var_embed}
\texttt{VarEmbed(default\_vars, embed\_dim)} injects a learnable identifier for each physical variable (e.g., temperature, pressure, etc.). Given \texttt{in\_vars}, the module adds/concatenates a variable-specific embedding to the corresponding token stream:
\[
z_1 = \texttt{VarEmbed}(z_0, \texttt{in\_vars}).
\]
This makes the representation explicitly aware of which variable generated each token.

\subsubsection{(4) Aggregation Across Variables via Attention}
\label{subsec:aggvars}
After variable tagging, \texttt{AggVars(embed\_dim, num\_heads)} aggregates the variable-wise tokenizations using an attention mechanism (multi-head attention controlled by \texttt{num\_heads}):
\[
z_2 = \texttt{AggVars}(z_1).
\]
Conceptually, this step fuses information across variables at the patch-token level, producing a joint latent representation that mixes multi-variable context.

\subsubsection{(5) Patch Position Encoding}
\label{subsec:pos_embed}
\texttt{PosEmbed(embed\_dim, patch\_size, image\_size, num\_patches)} encodes patch indices/positions:
\[
z_3 = \texttt{PosEmbed}(z_2).
\]
As stated in \texttt{bomberman.py}, the position encoding is ``only index-aware'' and does not necessarily encode continuous geometry; it primarily gives the transformer-style token sequence a notion of ordering.

\subsubsection{(6) Lead-Time (Temporal) Encoding}
\label{subsec:time_embed}
\texttt{TimeEmbed(embed\_dim)} injects a temporal offset conditioning (e.g., forecast lead time):
\[
z_4 = \texttt{TimeEmbed}(z_3, \texttt{lead\_times}).
\]
This enables the same network to represent different time horizons through a learned embedding of \texttt{lead\_times}.

\subsubsection{(7) SWIN U-Net Backbone on Patch Tokens}
\label{subsec:swin_unet}
The core spatial modeling is performed by \texttt{SwinUnetBackbone}:
\[
z_5 = \texttt{SwinUnetBackbone}(z_4).
\]
It is parameterized by:
\begin{itemize}
    \item base token width $D=$ \texttt{embed\_dim} and scale factor \texttt{emb\_factor} controlling feature expansion/contraction across U-Net stages,
    \item \texttt{patch\_grid\_size} derived from the patch embedding grid,
    \item \texttt{block\_structure} (e.g., \texttt{(1,1,3,1)}) specifying the number of SWIN encoder blocks per level,
    \item \texttt{window\_sizes} controlling local self-attention windows per stage (e.g., \texttt{(8,8)}, \texttt{(8,8)}, \texttt{(4,4)}, \texttt{(2,2)}),
    \item \texttt{patch\_merge\_scales} controlling down/up-sampling between levels (e.g., repeated \texttt{(2,2)}),
    \item \texttt{num\_heads} for multi-head self-attention.
\end{itemize}
This module implements a U-Net-like encoder--decoder with hierarchical patch merging/expanding and windowed self-attention (SWIN-style) for efficient high-resolution processing.

\subsubsection{(8) Projection to Per-Variable Patch Pixels}
\label{subsec:linear4unpatch}
After the SWIN U-Net, \texttt{LodeRunner} maps each token back to per-variable patch pixels using a linear layer:
\[
z_6 = W z_5 + b,\qquad
W \in \mathbb{R}^{(V P_h P_w)\times D}.
\]
In code this is \texttt{nn.Linear(embed\_dim, max\_vars * patch\_size[0] * patch\_size[1])} (\texttt{bomberman.py}). This produces, per patch token, the raw pixels for every modeled variable and every location inside the patch.

\subsubsection{(9) Unpatchify and Output Variable Selection}
\label{subsec:unpatchify}
Finally, \texttt{Unpatchify(total\_num\_vars=max\_vars, patch\_grid\_size, patch\_size)} reconstructs full-resolution variable fields:
\[
\hat{y} = \texttt{Unpatchify}(z_6)\in \mathbb{R}^{B\times V\times H\times W}.
\]
The module then returns only the requested outputs:
\[
\texttt{preds} = \hat{y}[:, \texttt{out\_vars}],
\]
matching the paper’s notion of a variable-conditioned multi-output surrogate model (\texttt{bomberman.py}).

\paragraph{Validity constraints.}
Before constructing submodules, \texttt{validate\_patch\_and\_window} checks that \texttt{image\_size}, \texttt{patch\_size}, \texttt{window\_sizes}, and \texttt{patch\_merge\_scales} are mutually compatible (e.g., divisibility across hierarchical stages). The model asserts all checks pass, enforcing consistent tensor shapes throughout the SWIN U-Net pipeline (\texttt{bomberman.py}).

\paragraph{Summary.}
Overall, \texttt{LodeRunner} combines (i) variable-wise patch tokenization, (ii) explicit variable identity encoding and cross-variable attention fusion, (iii) position and lead-time conditioning, and (iv) a hierarchical SWIN U-Net backbone, followed by a linear token-to-pixel projection and unpatchification to produce dense multi-variable image outputs.

\end{document}
